%%%%%%%%%%%%%%%%%%%%%%%%%%%%%%%%%%%%%%%%%
% Beamer Presentation
% LaTeX Template
% Version 2.0 (Oct, 10 2025)
%
% This template originates from:
% https://www.LaTeXTemplates.com
%
% Author:
% William Aisawa (alveswill@usp.br)
%
% License:
% CC BY-NC-SA 4.0 (https://creativecommons.org/licenses/by-nc-sa/4.0/)
%
%%%%%%%%%%%%%%%%%%%%%%%%%%%%%%%%%%%%%%%%%

%----------------------------------------------------------------------------------------
%	PACKAGES AND OTHER DOCUMENT CONFIGURATIONS
%----------------------------------------------------------------------------------------


\documentclass[
	11pt, % Set the default font size, options include: 8pt, 9pt, 10pt, 11pt, 12pt, 14pt, 17pt, 20pt
	%t, % Uncomment to vertically align all slide content to the top of the slide, rather than the default centered
	%aspectratio=169, % Uncomment to set the aspect ratio to a 16:9 ratio which matches the aspect ratio of 1080p and 4K screens and projectors
]{beamer}

\graphicspath{{images/}{./}} % Specifies where to look for included images (trailing slash required)

\usepackage{booktabs} % Allows the use of \toprule, \midrule and \bottomrule for better rules in tables
\usepackage{listings} % Para formatação de código (necessário para code_lang.tex)
\usepackage{xcolor}   % Para cores em código (necessário para code_lang.tex)

%----------------------------------------------------------------------------------------
%	SELECT LAYOUT THEME
%----------------------------------------------------------------------------------------

% Beamer comes with a number of default layout themes which change the colors and layouts of slides. Below is a list of all themes available, uncomment each in turn to see what they look like.

%\usetheme{default}
%\usetheme{AnnArbor}
%\usetheme{Antibes}
%\usetheme{Bergen}
%\usetheme{Berkeley}
%\usetheme{Berlin}
\usetheme{Boadilla}
%\usetheme{CambridgeUS}
%\usetheme{Copenhagen}
%\usetheme{Darmstadt}
%\usetheme{Dresden}
%\usetheme{Frankfurt}
%\usetheme{Goettingen}
%\usetheme{Hannover}
%\usetheme{Ilmenau}
%\usetheme{JuanLesPins}
%\usetheme{Luebeck}
%\usetheme{Madrid}
%\usetheme{Malmoe}
%\usetheme{Marburg}
%\usetheme{Montpellier}
%\usetheme{PaloAlto}
%\usetheme{Pittsburgh}
%\usetheme{Rochester}
%\usetheme{Singapore}
%\usetheme{Szeged}
%\usetheme{Warsaw}

%----------------------------------------------------------------------------------------
%	SELECT COLOR THEME
%----------------------------------------------------------------------------------------

% Beamer comes with a number of color themes that can be applied to any layout theme to change its colors. Uncomment each of these in turn to see how they change the colors of your selected layout theme.

%\usecolortheme{albatross}
%\usecolortheme{beaver}
%\usecolortheme{beetle}
%\usecolortheme{crane}
%\usecolortheme{dolphin}
%\usecolortheme{dove}
%\usecolortheme{fly}
%\usecolortheme{lily}
%\usecolortheme{monarca}
%\usecolortheme{seagull}
%\usecolortheme{seahorse}
%\usecolortheme{spruce}
%\usecolortheme{whale}
%\usecolortheme{wolverine}


%----------------------------------------------------------------------------------------
%	SELECT FONT THEME & FONTS
%----------------------------------------------------------------------------------------

% Beamer comes with several font themes to easily change the fonts used in various parts of the presentation. Review the comments beside each one to decide if you would like to use it. Note that additional options can be specified for several of these font themes, consult the beamer documentation for more information.

\usefonttheme{default} % Typeset using the default sans serif font
%\usefonttheme{serif} % Typeset using the default serif font (make sure a sans font isn't being set as the default font if you use this option!)
%\usefonttheme{structurebold} % Typeset important structure text (titles, headlines, footlines, sidebar, etc) in bold
%\usefonttheme{structureitalicserif} % Typeset important structure text (titles, headlines, footlines, sidebar, etc) in italic serif
%\usefonttheme{structuresmallcapsserif} % Typeset important structure text (titles, headlines, footlines, sidebar, etc) in small caps serif

%------------------------------------------------

%\usepackage{mathptmx} % Use the Times font for serif text
\usepackage{palatino} % Use the Palatino font for serif text
% \usepackage{booktabs} % (Removido, já está acima)
%\usepackage{helvet} % Use the Helvetica font for sans serif text
\usepackage[default]{opensans} % Use the Open Sans font for sans serif text
%\usepackage[default]{FiraSans} % Use the Fira Sans font for sans serif text
%\usepackage[default]{lato} % Use the Lato font for sans serif text
\usepackage{graphicx}
\usepackage{caption}
\usepackage{subcaption}
\usepackage{float} 



%----------------------------------------------------------------------------------------
%	SELECT INNER THEME
%----------------------------------------------------------------------------------------

% Inner themes change the styling of internal slide elements, for example: bullet points, blocks, bibliography entries, title pages, theorems, etc. Uncomment each theme in turn to see what changes it makes to your presentation.

%\useinnertheme{default}
\useinnertheme{circles}
%\useinnertheme{rectangles}
%\useinnertheme{rounded}
%\useinnertheme{inmargin}

%----------------------------------------------------------------------------------------
%	SELECT OUTER THEME
%----------------------------------------------------------------------------------------

% Outer themes change the overall layout of slides, such as: header and footer lines, sidebars and slide titles. Uncomment each theme in turn to see what changes it makes to your presentation.

%\useoutertheme{default}
%\useoutertheme{infolines}
%\useoutertheme{miniframes}
%\useoutertheme{smoothbars}
%\useoutertheme{sidebar}
%\useoutertheme{split}
%\useoutertheme{shadow}
%\useoutertheme{tree}
%\useoutertheme{smoothtree}

%\setbeamertemplate{footline} % Uncomment this line to remove the footer line in all slides
%\setbeamertemplate{footline}[page number] % Uncomment this line to replace the footer line in all slides with a simple slide count

\setbeamertemplate{navigation symbols}{} % Uncomment this line to remove the navigation symbols from the bottom of all slides

% Cores Personalizadas
\definecolor{primaryColor}{RGB}{20,45,105}   % Cor primária - azul escuro
\definecolor{secondaryColor}{RGB}{0,100,160} % Cor secundária - azul médio

% Configurações de Cores
\setbeamercolor{structure}{fg=primaryColor}
\setbeamercolor{palette primary}{bg=primaryColor, fg=white}
\setbeamercolor{palette secondary}{bg=secondaryColor, fg=white}
\setbeamercolor{title}{bg=primaryColor, fg=white}

% Cores do Cabeçalho e Rodapé
\setbeamercolor{headline}{bg=secondaryColor, fg=white}
\setbeamercolor{section in head/foot}{bg=primaryColor, fg=white}
\setbeamercolor{subsection in head/foot}{bg=secondaryColor, fg=white}
\setbeamercolor{author in head/foot}{bg=primaryColor, fg=white}
\setbeamercolor{title in head/foot}{bg=secondaryColor, fg=white}
\setbeamercolor{date in head/foot}{bg=primaryColor, fg=white}
\setbeamercolor{page number in head/foot}{bg=primaryColor, fg=white}

%----------------------------------------------------------------------------------------
% BIBLIOGRAFIA
%----------------------------------------------------------------------------------------
% Removido biblatex/biber para compatibilidade com Overleaf gratuito
% Utilize os comandos tradicionais do BibTeX no final do documento:
%   \bibliographystyle{alpha}
%   \bibliography{bibliography}
%----------------------------------------------------------------------------------------
%	PRESENTATION INFORMATION
%----------------------------------------------------------------------------------------

\title[Short Title]{Full Presentation Title} % The short title in the optional parameter appears at the bottom of every slide, the full title in the main parameter is only on the title page

\subtitle{Optional Subtitle} % Presentation subtitle, remove this command if a subtitle isn't required

\author[William Akihiro Alves Aisawa \and Sarita Mazzini Bruschi]{William Aisawa \and Sarita Bruschi} % Presenter name(s), the optional parameter can contain a shortened version to appear on the bottom of every slide, while the main parameter will appear on the title slide

\institute[UC]{Universidade de São Paulo  \\ \smallskip \textit{alveswill@usp.br}} % Your institution, the optional parameter can be used for the institution shorthand and will appear on the bottom of every slide after author names, while the required parameter is used on the title slide and can include your email address or additional information on separate lines

\date[\today]{International Symposium of Explorers \\ \today} % Presentation date or conference/meeting name, the optional parameter can contain a shortened version to appear on the bottom of every slide, while the required parameter value is output to the title slide


%----------------------------------------------------------------------------------------
%	BEGIN PRESENTATION
%----------------------------------------------------------------------------------------
\begin{document}
%----------------------------------------------------------------------------------------
% CONFIGURAÇÕES PARA CÓDIGO
%----------------------------------------------------------------------------------------


% Cores para syntax highlighting (VSCode Light Theme)
\definecolor{vscBackground}{RGB}{255,255,255}    % Fundo branco
\definecolor{vscKeyword}{RGB}{175,0,219}         % Roxo para palavras-chave
\definecolor{vscString}{RGB}{163,21,21}          % Vermelho para strings
\definecolor{vscComment}{RGB}{0,128,0}           % Verde para comentários
\definecolor{vscFunction}{RGB}{121,94,38}        % Marrom para funções
\definecolor{vscNumber}{RGB}{9,134,88}           % Verde escuro para números
\definecolor{vscOperator}{RGB}{175,0,219}        % Roxo para operadores
\definecolor{vscText}{RGB}{0,0,0}                % Texto preto
\definecolor{vscLineNr}{RGB}{128,128,128}        % Cinza para números de linha

% Configuração geral do listings para UTF-8
\lstset{
    inputencoding=utf8,
    extendedchars=true,
    literate=%
        {á}{{\'a}}1 {é}{{\'e}}1 {í}{{\'i}}1 {ó}{{\'o}}1 {ú}{{\'u}}1
        {Á}{{\'A}}1 {É}{{\'E}}1 {Í}{{\'I}}1 {Ó}{{\'O}}1 {Ú}{{\'U}}1
        {à}{{\`a}}1 {è}{{\`e}}1 {ì}{{\`i}}1 {ò}{{\`o}}1 {ù}{{\`u}}1
        {À}{{\`A}}1 {È}{{\'E}}1 {Ì}{{\`I}}1 {Ò}{{\`O}}1 {Ù}{{\`U}}1
        {ã}{{\~a}}1 {ẽ}{{\~e}}1 {ĩ}{{\~i}}1 {õ}{{\~o}}1 {ũ}{{\~u}}1
        {Ã}{{\~A}}1 {Ẽ}{{\~E}}1 {Ĩ}{{\~I}}1 {Õ}{{\~O}}1 {Ũ}{{\~U}}1
        {â}{{\^a}}1 {ê}{{\^e}}1 {î}{{\^i}}1 {ô}{{\^o}}1 {û}{{\^u}}1
        {Â}{{\^A}}1 {Ê}{{\^E}}1 {Î}{{\^I}}1 {Ô}{{\^O}}1 {Û}{{\^U}}1
        {ç}{{\c c}}1 {Ç}{{\c C}}1
        {º}{{\textordmasculine}}1
        {ª}{{\textordfeminine}}1
}

% Configurações base comum para todas as linguagens
\lstdefinestyle{baseStyle}{
    backgroundcolor=\color{vscBackground},
    basicstyle=\ttfamily\small\color{vscText},
    breakatwhitespace=false,
    breaklines=true,
    captionpos=b,
    keepspaces=true,
    numbers=left,
    numbersep=5pt,
    showspaces=false,
    showstringspaces=false,
    showtabs=false,
    tabsize=4,
    frame=single,
    framerule=0.8pt,
    rulecolor=\color{gray!20},
    numberstyle=\tiny\color{vscLineNr},
    keywordstyle=\color{vscKeyword},
    commentstyle=\color{vscComment}\itshape,
    stringstyle=\color{vscString},
    emphstyle=\color{vscFunction},
    columns=flexible,
    basewidth={0.5em,0.45em},
    inputencoding=utf8,
    extendedchars=true
}

%----------------------------------------------------------------------------------------
% Python
%----------------------------------------------------------------------------------------
\lstdefinestyle{pythonStyle}{
    style=baseStyle,
    language=Python,
    morekeywords={self,None,True,False,import,from,as,def,class,return,yield,
                  for,while,if,else,elif,try,except,finally,with,lambda,
                  async,await,break,continue,global,nonlocal,pass,raise},
    morekeywords=[2]{print,len,range,type,int,str,float,list,dict,set,
                     tuple,max,min,sum,sorted,enumerate,zip,map,filter,
                     any,all,abs,round,pow,divmod},
    keywordstyle=[2]\color{vscFunction},
    sensitive=true
}

\lstnewenvironment{python}[1][]{\lstset{style=pythonStyle, #1}}{}
\newcommand{\pyinline}[1]{\lstinline[style=pythonStyle]!#1!}
\newcommand{\inputpython}[2][]{\lstinputlisting[style=pythonStyle,#1]{#2}}

%----------------------------------------------------------------------------------------
% C Language
%----------------------------------------------------------------------------------------
\lstdefinestyle{cStyle}{
    style=baseStyle,
    language=C,
    morekeywords={include,define,void,int,char,float,double,long,unsigned,
                  struct,union,enum,typedef,const,static,extern,register,
                  auto,volatile,sizeof,return,if,else,for,while,do,switch,
                  case,break,continue,default,goto},
    morekeywords=[2]{printf,scanf,malloc,free,calloc,realloc,fopen,fclose,
                     fprintf,fscanf,strcpy,strlen,strcat},
    keywordstyle=[2]\color{vscFunction},
    sensitive=true
}

\lstnewenvironment{clang}[1][]{\lstset{style=cStyle, #1}}{}
\newcommand{\clinline}[1]{\lstinline[style=cStyle]!#1!}
\newcommand{\inputclang}[2][]{\lstinputlisting[style=cStyle,#1]{#2}}

%----------------------------------------------------------------------------------------
% C++
%----------------------------------------------------------------------------------------
\lstdefinestyle{cppStyle}{
    style=baseStyle,
    language=C++,
    morekeywords={class,private,protected,public,template,typename,namespace,
                  using,new,delete,this,friend,virtual,override,final,explicit,
                  mutable,constexpr,nullptr,noexcept,static_cast,dynamic_cast,
                  const_cast},
    morekeywords=[2]{cout,cin,endl,vector,string,map,set,queue,stack,pair,
                     begin,end,push_back,pop_back,emplace_back,size,empty},
    keywordstyle=[2]\color{vscFunction},
    sensitive=true
}

\lstnewenvironment{cpp}[1][]{\lstset{style=cppStyle, #1}}{}
\newcommand{\cppinline}[1]{\lstinline[style=cppStyle]!#1!}
\newcommand{\inputcpp}[2][]{\lstinputlisting[style=cppStyle,#1]{#2}}

%----------------------------------------------------------------------------------------
% R Language
%----------------------------------------------------------------------------------------
\lstdefinestyle{rStyle}{
    style=baseStyle,
    language=R,
    morekeywords={if,else,repeat,while,function,for,in,next,break,TRUE,FALSE,
                  NULL,Inf,NaN,NA,NA_integer_,NA_real_,NA_complex_,NA_character_},
    morekeywords=[2]{library,require,attach,detach,source,setwd,options,
                     data.frame,read.csv,write.csv,list,matrix,array},
    keywordstyle=[2]\color{vscFunction},
    sensitive=true
}

\lstnewenvironment{rlang}[1][]{\lstset{style=rStyle, #1}}{}
\newcommand{\rlinline}[1]{\lstinline[style=rStyle]!#1!}
\newcommand{\inputrlang}[2][]{\lstinputlisting[style=rStyle,#1]{#2}}

%----------------------------------------------------------------------------------------
% Java
%----------------------------------------------------------------------------------------
\lstdefinestyle{javaStyle}{
    style=baseStyle,
    language=Java,
    morekeywords={abstract,assert,boolean,break,byte,case,catch,char,class,
                  const,continue,default,do,double,else,enum,extends,final,
                  finally,float,for,if,implements,import,instanceof,int,
                  interface,long,native,new,package,private,protected,public,
                  return,short,static,strictfp,super,switch,synchronized,this,
                  throw,throws,transient,try,void,volatile,while},
    morekeywords=[2]{String,System,out,println,printStackTrace,ArrayList,
                     HashMap,Arrays,List,Map,Set,Exception,RuntimeException},
    keywordstyle=[2]\color{vscFunction},
    sensitive=true
}

\lstnewenvironment{java}[1][]{\lstset{style=javaStyle, #1}}{}
\newcommand{\javainline}[1]{\lstinline[style=javaStyle]!#1!}
\newcommand{\inputjava}[2][]{\lstinputlisting[style=javaStyle,#1]{#2}}
\begin{frame}
%----------------------------------------------------------------------------------------
%	LOGOS
%----------------------------------------------------------------------------------------
    \centering
    \begin{minipage}{0.25\textwidth}
        \centering
        \includegraphics[width=0.7\linewidth]{images/lab_logo.jpeg}
    \end{minipage}%
    \hfill
    \begin{minipage}{0.25\textwidth}
        \centering
        \includegraphics[width=0.55\linewidth]{images/icmc_logo.png}
    \end{minipage}
%----------------------------------------------------------------------------------------
%	TITLE PAGE
%----------------------------------------------------------------------------------------

    \titlepage 
\end{frame}

%----------------------------------------------------------------------------------------
%	BEGIN SUMMARY
%----------------------------------------------------------------------------------------

\begin{frame}
    \frametitle{Estrutura da apresentação}
    \tableofcontents
\end{frame}

%----------------------------------------------------------------------------------------
%	INCLUDE SECTIONS
%----------------------------------------------------------------------------------------

\section{Exemplos com texto} % Seções são adicionadas para organizar sua apresentação em blocos discretos, todas as seções e subseções são automaticamente exibidas no índice como uma visão geral da apresentação, mas NÃO são exibidas como slides separados.

%----------------------------------------------------------------------------------------

\begin{frame}
	\frametitle{Texto corrido}
    Lorem ipsum dolor sit amet, consectetur adipiscing elit. Nullam ipsum velit, cursus quis ligula eu, malesuada aliquet massa. Quisque non convallis felis, a auctor eros. Etiam sit amet turpis a sapien pulvinar malesuada quis quis nisi. Quisque scelerisque volutpat ligula vel mollis. Nam sit amet tristique erat, sit amet cursus mi. 
\end{frame}

%----------------------------------------------------------------------------------------

\begin{frame}
	\frametitle{Texto em tópicos numerados}
     Lorem ipsum dolor sit amet, consectetur adipiscing elit:
    \begin{enumerate}
        \item Lorem ipsum dolor sit amet.
        \item Lorem ipsum dolor sit amet.
    \end{enumerate}
	
\end{frame}

%----------------------------------------------------------------------------------------

\begin{frame}
	\frametitle{Texto em tópicos}
     Lorem ipsum dolor sit amet, consectetur adipiscing elit:
    \begin{itemize}
        \item Lorem ipsum dolor sit amet.
        \item Lorem ipsum dolor sit amet.
    \end{itemize}
	
\end{frame}



\section{Exemplos com equações e imagens} % Seções são adicionadas para organizar sua apresentação em blocos discretos, todas as seções e subseções são automaticamente exibidas no índice como uma visão geral da apresentação, mas NÃO são exibidas como slides separados.

%----------------------------------------------------------------------------------------

\begin{frame}{Uma imagem}
   \begin{figure}[h]
       \centering
       \includegraphics[width=0.7\textwidth]{images/lab_logo.jpeg}
       \caption{Legenda da imagem}
       \label{fig:label_da_imagem}
   \end{figure}
\end{frame}

%----------------------------------------------------------------------------------------

\begin{frame}{Duas imagens}
   \begin{figure}
       \centering
       \begin{subfigure}[b]{0.45\textwidth}
           \centering
           \includegraphics[width=\textwidth]{images/lab_logo.jpeg}
           \caption{Legenda 1}
           \label{fig:img1}
       \end{subfigure}
       \hfill
       \begin{subfigure}[b]{0.45\textwidth}
           \centering
           \includegraphics[width=\textwidth]{images/lab_logo.jpeg}
           \caption{Legenda 2}
           \label{fig:img2}
       \end{subfigure}
   \end{figure}
\end{frame}

%----------------------------------------------------------------------------------------
\begin{frame}{Equações}
    Equações de Navier-Stokes Forma expandida (3D):
    \footnotesize
        \begin{align*}
            \rho\left(\frac{\partial u}{\partial t} + u\frac{\partial u}{\partial x} + v\frac{\partial u}{\partial y} + w\frac{\partial u}{\partial z}\right) &= -\frac{\partial p}{\partial x} + \mu\left(\frac{\partial^2 u}{\partial x^2} + \frac{\partial^2 u}{\partial y^2} + \frac{\partial^2 u}{\partial z^2}\right) + f_x \\[0.3cm]
            \rho\left(\frac{\partial v}{\partial t} + u\frac{\partial v}{\partial x} + v\frac{\partial v}{\partial y} + w\frac{\partial v}{\partial z}\right) &= -\frac{\partial p}{\partial y} + \mu\left(\frac{\partial^2 v}{\partial x^2} + \frac{\partial^2 v}{\partial y^2} + \frac{\partial^2 v}{\partial z^2}\right) + f_y \\[0.3cm]
            \rho\left(\frac{\partial w}{\partial t} + u\frac{\partial w}{\partial x} + v\frac{\partial w}{\partial y} + w\frac{\partial w}{\partial z}\right) &= -\frac{\partial p}{\partial z} + \mu\left(\frac{\partial^2 w}{\partial x^2} + \frac{\partial^2 w}{\partial y^2} + \frac{\partial^2 w}{\partial z^2}\right) + f_z
        \end{align*}
        
    onde $\mathbf{v} = (u,v,w)$ é o campo de velocidade, $p$ é a pressão, $\rho$ é a densidade, $\mu$ é a viscosidade dinâmica e $\mathbf{f}$ representa forças externas.
\end{frame}
\section{Exemplos com código} % Seções são adicionadas para organizar sua apresentação em blocos discretos, todas as seções e subseções são automaticamente exibidas no índice como uma visão geral da apresentação, mas NÃO são exibidas como slides separados.

%----------------------------------------------------------------------------------------
\begin{frame}[fragile]
    \frametitle{Python}
    
    \begin{python}
def calcular_dobro(x):
    """Retorna o dobro do número"""
    return 2 * x

# Testando a função
numero = 5
resultado = calcular_dobro(numero)
print(f"O dobro de {numero} é {resultado}")
    \end{python}
\end{frame}

%----------------------------------------------------------------------------------------
\begin{frame}[fragile]
    \frametitle{C}
    
    \begin{clang}
#include <stdio.h>

int main() {
    int numero = 5;
    int dobro = 2 * numero;
    
    printf("O dobro de %d eh %d\n", numero, dobro);
    return 0;
}
    \end{clang}
\end{frame}

%----------------------------------------------------------------------------------------
\begin{frame}[fragile]
    \frametitle{C++}
    
    \begin{cpp}
#include <iostream>
using namespace std;

int main() {
    int numero = 5;
    int dobro = 2 * numero;
    
    cout << "O dobro de " << numero;
    cout << " eh " << dobro << endl;
    return 0;
}
    \end{cpp}
\end{frame}

%----------------------------------------------------------------------------------------
\begin{frame}[fragile]
    \frametitle{R}
    
    \begin{rlang}
# Função para calcular o dobro
calcular_dobro <- function(x) {
  return(2 * x)
}

# Testando a função
numero <- 5
resultado <- calcular_dobro(numero)
print(paste("O dobro de", numero, "é", resultado))
    \end{rlang}
\end{frame}

%----------------------------------------------------------------------------------------

\begin{frame}[fragile]
    \frametitle{Java}
    
    \begin{java}
public class Exemplo {
    public static void main(String[] args) {
        int numero = 5;
        int dobro = 2 * numero;
        
        System.out.println("O dobro de " + numero +
                         " eh " + dobro);
    }
}
    \end{java}
\end{frame}
\section{Conclusão} % Seções são adicionadas para organizar sua apresentação em blocos discretos, todas as seções e subseções são automaticamente exibidas no índice como uma visão geral da apresentação, mas NÃO são exibidas como slides separados.




%----------------------------------------------------------------------------------------
%	FINAL SLIDE
%----------------------------------------------------------------------------------------
\begin{frame}
    \begin{center}
        {\Huge Fim da apresentação!}
    \end{center}
\end{frame}

%----------------------------------------------------------------------------------------
% REFERÊNCIAS
%----------------------------------------------------------------------------------------

\begin{frame}{Referências}
    \nocite{*}
    % \printbibliography[heading=none] % (Removido para uso com BibTeX)

\bibliographystyle{alpha}
\bibliography{bibliography}
\end{frame}
\end{document}